\section{Conclusion}

Language models show recurring teaching tendencies, but the evidence differs
across behaviors. Helping choices have the strongest combined support; other
measures show differences within narrower settings, shared responses, or mixed
results. Knowing how often a model gives answers or invites reasoning helps
predict these actions on new problems, even under tested changes in student
wording. For acknowledgement, knowing how each model responds to student cues
adds predictive value beyond its usual rate on new problems, but this extra
benefit is lost under new wording. Explicit instructions can make models perform
the same required action while leaving other choices different.

These findings help educators compare tutors through their teaching habits
alongside task capability. They also help developers investigate which gaps
between a model's behavior and an activity's requirements need better prompts
or additional system support.
